\documentclass[12pt,a4paper]{article}

\usepackage[margin=2.5cm]{geometry}
\usepackage{amsmath,amssymb,amsthm}


\usepackage{booktabs}
\usepackage{graphicx}
\usepackage[colorlinks=true,linkcolor=blue,citecolor=blue,urlcolor=blue]{hyperref}
\usepackage[round]{natbib}
\usepackage{xcolor}
\usepackage{microtype}
\usepackage{lmodern}
\usepackage[T1]{fontenc}
\usepackage{setspace}
\onehalfspacing

\newcommand{\pkg}[1]{\textbf{#1}}
\newcommand{\code}[1]{\texttt{#1}}
\newcommand{\email}[1]{\texttt{#1}}

\newtheorem{definition}{Definition}
\newtheorem{proposition}{Proposition}
\newtheorem{remark}{Remark}

\title{\textbf{\pkg{lepidoptera}: A Domain-Agnostic Engine for Causal Simulation\\
with Mixed Variable Types, Vine Copula Dependence,\\
and Counterfactual Propagation}\\[1em]
\large \textit{Draft — not for citation}}

\author{Danilo A. Sarti\\
\small \email{daniloasarti@gmail.com}}

\date{\today}

\begin{document}
\maketitle
\tableofcontents
\newpage

%% ─── Abstract ────────────────────────────────────────────────────────────────
\begin{abstract}
We introduce \pkg{lepidoptera}, an open-source R and Python package providing a modular, domain-agnostic engine for structural causal model (SCM) simulation with mixed variable types. Unlike existing tools that restrict mechanisms to continuous variables or treat interventions as post-hoc operations, \pkg{lepidoptera} elevates the mechanism---the generating rule associated with each node in a directed acyclic graph (DAG)---to a first-class, inspectable, and serialisable object. The framework natively supports all nine type-pair combinations (continuous, categorical, and ordinal in any parent-child configuration), a library of pre-built mechanisms including multinomial logit, ordered logit, group-effect, and categorical transition models, and a vine copula mechanism that learns arbitrary dependence structures from observed data via the Rosenblatt transform. A key design feature is \emph{partial propagation}: by freezing the random draws of unaffected nodes and replaying only the downstream consequences of an intervention, \pkg{lepidoptera} enables individual-level counterfactual reasoning without additional assumptions beyond the original SCM. An AI-native extension layer allows any pre-fitted predictive model---random forests, gradient-boosted trees, neural networks---to serve as a generative mechanism through a single \code{predict\_fn} interface. Comprehensive simulation studies validate type-combination coverage, copula family fidelity, and counterfactual consistency. Real-data illustrations span clinical trial planning, hospital emergency-ward digital twinning, and epidemiological mediation analysis. \pkg{lepidoptera} is available under the MIT licence.

\bigskip\noindent\textbf{Keywords:} structural causal model, directed acyclic graph, mixed variable types, vine copula, counterfactual simulation, partial propagation, R, Python
\end{abstract}

%% ─── 1. Introduction ─────────────────────────────────────────────────────────
\section{Introduction}
\label{sec:intro}

Simulation from structural causal models (SCMs) has become an essential tool across empirical disciplines. Epidemiologists use it to plan randomised trials under realistic confounding \citep{hernan2020causal}; economists employ it to evaluate counterfactual policy scenarios \citep{angrist2009mostly}; clinical data scientists use it to stress-test machine-learning pipelines before deployment \citep{vanderschueren2023}; and computer scientists increasingly rely on synthetic data with known causal structure to benchmark causal discovery algorithms \citep{spirtes2000causation, glymour2019review}.

Despite this broad demand, the software landscape remains fragmented. Available tools either (i) restrict generating mechanisms to continuous Gaussian variables \citep{scutari2010bnlearn}, (ii) focus on DAG visualisation without generative capability \citep{textor2016dagitty}, (iii) hard-wire longitudinal study designs \citep{sofrygin2017simcausal}, or (iv) conflate the causal structure with a specific model family such as structural equation models \citep{holst2013lava, rosseel2012lavaan}. No existing package provides a unified, domain-agnostic engine in which the mechanism at every node is an explicit, replaceable object, interventions are first-class operations on that object graph, all three variable types are supported without workarounds, and individual-level (counterfactual) and population-level (interventional) queries are cleanly separated.

We present \pkg{lepidoptera}, an open-source package available in R and Python with a shared conceptual API, that fills this gap. The name is a nod to the systematic, modular beauty of biological taxonomy: just as a lepidopterist classifies organisms by explicit, composable features, the package classifies data-generating processes by their constituent mechanisms, each typed and independently swappable.

The central architectural contribution is the separation of \emph{structure} (which variables exist and who depends on whom, encoded in a DAG) from \emph{mechanism} (how each variable is generated given its parents, encoded in an explicit \code{Mechanism} object). This separation, grounded in the modularity property of SCMs \citep{pearl2009causality, peters2017elements}, has three practical consequences. First, any node mechanism can be replaced for any intervention or sensitivity analysis without re-specifying the rest of the system. Second, \emph{partial propagation}---recomputing only the downstream subgraph of an intervention while preserving the random draws of all unaffected nodes---is exact and requires no additional assumptions, yielding individual-level counterfactual queries in the sense of \citet{pearl2009causality}. Third, the mechanism slot can hold any callable respecting a minimal signature, so pre-fitted machine-learning models plug in directly via \code{model\_mechanism()}.

The paper is organised as follows. Section~\ref{sec:related} reviews related work. Section~\ref{sec:framework} gives formal definitions. Section~\ref{sec:mechanisms} describes the mechanism library. Section~\ref{sec:simulation} covers the simulation and propagation algorithms. Section~\ref{sec:counterfactual} formalises counterfactual propagation. Section~\ref{sec:copula} presents the vine copula mechanism. Section~\ref{sec:ai} describes the AI-native extension layer. Section~\ref{sec:sims} reports simulation studies. Section~\ref{sec:applications} presents real-data applications. Section~\ref{sec:discussion} concludes.

%% ─── 2. Related Work ─────────────────────────────────────────────────────────
\section{Related Work}
\label{sec:related}

\subsection{Causal simulation tools}

\citet{sofrygin2017simcausal} introduced \pkg{simcausal} for R, providing a rich language for specifying longitudinal observational and randomised-study designs with time-varying treatments and censoring. The package has been widely used in causal inference methodology papers \citep{luedtke2016, laan2018} and is the closest predecessor to \pkg{lepidoptera}. However, \pkg{simcausal} is designed around a temporal index and a shared-environment expression evaluator. This makes it natural for longitudinal settings but unwieldy for cross-sectional DAGs with mixed variable types and heterogeneous mechanisms. There is no concept of a reusable, typed mechanism object, and counterfactual (same-individual) queries are not natively supported.

\citet{goldfeld2020simstudy} provide \pkg{simstudy}, an expressive R package for simulating study populations under a variety of distributional assumptions. Its formula-based interface is convenient for power calculations and trial design, but the package does not expose a causal graph layer and intervention propagation must be managed manually.

\pkg{bnlearn} \citep{scutari2010bnlearn} focuses on Bayesian network structure learning and inference. Mechanisms are restricted to linear Gaussian or discrete conditional probability tables; there is no support for ordinal variables with proportional-odds structure, vine copula dependence, or partial propagation.

For Python, \pkg{causalnex} \citep{beaumont2021causalnex} and \pkg{dowhy} \citep{sharma2020dowhy} provide causal inference pipelines, but neither is designed as a simulation engine. \pkg{causalnex} focuses on structure learning and Bayesian network inference, while \pkg{dowhy} focuses on identification and estimation from observational data.

\subsection{Structural causal models and counterfactuals}

The formal foundations draw on Pearl's SCM theory \citep{pearl2009causality}, in which a system is defined by structural equations $X_i = f_i(\mathbf{Pa}_i, U_i)$ with exogenous noise terms $U_i$. The do-operator $P(Y \mid do(X=x))$ captures population-level interventions \citep{pearl1995causal}. Individual-level counterfactuals $Y_{X=x}(u)$ require the three-step procedure of abduction, action, and prediction \citep{pearl2009causality}. The potential outcomes framework of \citet{rubin1974estimating} and \citet{neyman1923} provides complementary formalisation. SUTVA \citep{rubin1980randomization} is satisfied by construction in \pkg{lepidoptera}'s closed-world model.

\subsection{Mixed variable types in graphical models}

Mixed graphical models combining continuous and discrete variables have a long history \citep{lauritzen1989graphical}. In the context of directed models, practical simulation tools for the user-specified mixed-type DAG setting---which requires type-dispatch for every parent-type/child-type pair---have received little software attention. Table~\ref{tab:comparison} summarises the state of the art.

\begin{table}[ht]
\centering
\caption{Feature comparison of causal simulation tools. ``Mixed'' = native support for continuous, categorical, and ordinal. ``Copula'' = vine copula mechanism. ``Part.\ prop.'' = individual-level counterfactual via partial propagation. ``R/Py'' = availability in both languages.}
\label{tab:comparison}
\small
\begin{tabular}{lccccc}
\toprule
Package & Mixed & Copula & Part.\ prop. & Typed mech. & R/Py \\
\midrule
\pkg{simcausal} \citep{sofrygin2017simcausal}   & No  & No  & No  & No  & R \\
\pkg{simstudy}  \citep{goldfeld2020simstudy}    & Partial & No & No & No & R \\
\pkg{bnlearn}   \citep{scutari2010bnlearn}       & Partial & No & No & No & R \\
\pkg{dagitty}   \citep{textor2016dagitty}        & No (viz) & No & No & No & R \\
\pkg{lavaan}    \citep{rosseel2012lavaan}         & No  & No  & No  & No  & R \\
\pkg{causalnex} \citep{beaumont2021causalnex}   & Partial & No & No & No & Py \\
\pkg{dowhy}     \citep{sharma2020dowhy}          & No (estim.) & No & No & No & Py \\
\pkg{lepidoptera} (this paper)                  & \textbf{Yes} & \textbf{Yes} & \textbf{Yes} & \textbf{Yes} & R+Py \\
\bottomrule
\end{tabular}
\end{table}

%% ─── 3. Framework ────────────────────────────────────────────────────────────
\section{The \pkg{lepidoptera} Framework}
\label{sec:framework}

\subsection{Formal definitions}

Let $\mathcal{V} = \{V_1, \ldots, V_p\}$ be a finite set of variables. Each $V_i$ has a \emph{type} $\tau(V_i) \in \{\text{continuous}, \text{categorical}, \text{ordinal}\}$ and a \emph{domain} $\mathcal{D}(V_i)$.

\begin{definition}[Mechanism]
A \emph{mechanism} for $V_i$ with parent set $\mathbf{Pa}_i$ is a triple $\mathcal{M}_i = (\mathrm{fn}_i,\, \boldsymbol{\tau}_{\mathbf{Pa}},\, \tau_i)$ where $\mathrm{fn}_i\colon \prod_{j\in\mathbf{Pa}_i}\mathcal{D}(V_j)\times\Omega \to \mathcal{D}(V_i)$ is measurable, $\boldsymbol{\tau}_{\mathbf{Pa}}$ is the declared parent type vector, and $\tau_i$ is the declared output type.
\end{definition}

\begin{definition}[System]
A \emph{system} $\mathcal{S} = (\mathcal{V}, \mathcal{E}, \{\mathcal{M}_i\})$ is a DAG $(\mathcal{V},\mathcal{E})$ together with a mechanism $\mathcal{M}_i$ for each $V_i$, satisfying: (1) $G$ is acyclic; (2) every variable has exactly one mechanism; (3) declared parent types match actual feature types.
\end{definition}

\begin{definition}[Intervention and downstream closure]
An \emph{intervention} $\mathcal{I}$ modifies one or more mechanisms in one of four ways: \emph{do-fix} (constant), \emph{do-distribution} (new mechanism), \emph{shift} (transform on the output), or \emph{mechanism-override} (arbitrary replacement). Interventions are immutable and applied to a copy $\mathcal{S}'$ of $\mathcal{S}$. The \emph{affected set} is $\mathrm{Aff}(\mathcal{T},G) = \{V : \exists\text{ directed path from }\mathcal{T}\text{ to }V\text{ in }G\}$.
\end{definition}

%% ─── 4. Mechanism Library ────────────────────────────────────────────────────
\section{Mechanism Library}
\label{sec:mechanisms}

For each of the nine parent-type/child-type combinations, \pkg{lepidoptera} ships a default mechanism factory (Table~\ref{tab:mechanisms}).

\begin{table}[ht]
\centering
\caption{Pre-built mechanism factories. C = continuous, K = categorical, O = ordinal.}
\label{tab:mechanisms}
\small
\begin{tabular}{llp{8cm}}
\toprule
Parent & Child & Default mechanism \\
\midrule
C & C & \code{linear\_gaussian}: $Y = \alpha + \boldsymbol{\beta}^\top \mathbf{X} + \epsilon$, $\epsilon\sim\mathcal{N}(0,\sigma^2)$ \\
C & K & \code{multinomial\_logit}: $P(Y\!=\!k\mid\mathbf{X})\propto\exp(\alpha_k+\boldsymbol{\beta}_k^\top\mathbf{X})$ \\
C & O & \code{ordered\_logit}: proportional-odds model \citep{mccullagh1980regression} \\
K & C & \code{group\_effect}: $Y\mid X\!=\!k\sim\mathcal{N}(\mu_k,\sigma_k^2)$ \\
K & K & \code{categorical\_transition}: full conditional probability table \\
K & O & \code{categorical\_transition} with ordinal output \\
O & C & \code{group\_effect} or \code{linear\_gaussian} on numeric code \\
O & K & \code{categorical\_transition} \\
O & O & \code{categorical\_transition} with optional monotone constraint \\
\bottomrule
\end{tabular}
\end{table}

%% ─── 5. Simulation ───────────────────────────────────────────────────────────
\section{Simulation and Propagation}
\label{sec:simulation}

\paragraph{Algorithm 1: Simulate population.}
\begin{enumerate}
  \item Compute topological order $\pi$ of $G$ (Kahn's algorithm).
  \item For $i = 1,\ldots,n$: set seed $s+i$; for $k=1,\ldots,p$ in order: $V_{\pi_k} \leftarrow \mathrm{fn}_{\pi_k}\!\bigl(\{P_{i,\pi_j} : V_{\pi_j}\in\mathbf{Pa}_{\pi_k}\}\bigr)$.
\end{enumerate}

The per-unit seed $s+i$ ensures exact reproducibility and enables partial propagation. \code{finalize\_system()} runs five invariant checks (acyclicity, unique mechanisms, valid references, type consistency, reachability) before any simulation.

%% ─── 6. Counterfactual ───────────────────────────────────────────────────────
\section{Counterfactual Propagation}
\label{sec:counterfactual}

\paragraph{Algorithm 2: Partial propagation.}
\begin{enumerate}
  \item Compute $\mathrm{Aff}$; build modified system $\mathcal{S}'$.
  \item For $i=1,\ldots,n$: set seed $s+i$ (same as baseline).
  \item For each $V_{\pi_k}$ in topological order:
    \begin{itemize}
      \item If $V_{\pi_k}\notin\mathrm{Aff}$: $P^{\text{cf}}_{i,\pi_k} \leftarrow P^{\text{base}}_{i,\pi_k}$ (freeze).
      \item Else: advance RNG; $P^{\text{cf}}_{i,\pi_k} \leftarrow \mathrm{fn}'_{\pi_k}(\cdot)$.
    \end{itemize}
\end{enumerate}

\begin{proposition}[Counterfactual consistency]
If $V\notin\mathrm{Aff}(\mathcal{T},G)$ then $P^{\mathrm{cf}}_{i,V} = P^{\mathrm{base}}_{i,V}$ for all $i$. If $V\in\mathrm{Aff}(\mathcal{T},G)$ then $P^{\mathrm{cf}}_{i,V}$ is generated by the modified mechanism $\mathcal{M}'_V$ applied to the counterfactual parent values.
\end{proposition}

Table~\ref{tab:modes} contrasts the two modes.

\begin{table}[ht]
\centering
\caption{Population (\code{mode = "full"}) vs.\ individual-level (\code{mode = "partial"}) propagation.}
\label{tab:modes}
\small
\begin{tabular}{lll}
\toprule
Property & \code{mode = "full"} & \code{mode = "partial"} \\
\midrule
RNG streams & Re-drawn from scratch & Frozen for unaffected nodes \\
Estimand & ATE (policy) & ITE (counterfactual) \\
Use case & Trial design, policy & Explanation, fairness \\
\bottomrule
\end{tabular}
\end{table}

%% ─── 7. Vine Copula ──────────────────────────────────────────────────────────
\section{Vine Copula Mechanism}
\label{sec:copula}

When the analyst has observed data on a group of variables and wishes to preserve their joint dependence structure, \code{copula\_mechanism()} estimates a D-vine pair-copula construction \citep{aas2009pair, dissmann2013selecting} via \pkg{rvinecopulib} \citep{nagler2021rvinecopulib} and uses the Rosenblatt transform \citep{rosenblatt1952remarks} for conditional simulation.

Variables are ordered $V_1,\ldots,V_{d-1}, V_d$ with the \emph{child always last}. The D-vine structure ensures $v_d = F(u_d \mid u_1,\ldots,u_{d-1})$ under the Rosenblatt map $T$, which can be exactly inverted.

\paragraph{Algorithm 3: Copula conditional simulation (D-vine).}
\begin{enumerate}
  \item Set $\mathbf{u} \leftarrow (\hat{u}_1,\ldots,\hat{u}_{d-1}, 0.5)$ (placeholder for child column).
  \item $\mathbf{v} \leftarrow T(\mathbf{u};\hat{C})$ \quad (Rosenblatt transform).
  \item Replace $v_d \leftarrow w$ where $w\sim\mathrm{Uniform}(0,1)$.
  \item $\mathbf{u}^* \leftarrow T^{-1}(\mathbf{v};\hat{C})$ \quad (inverse Rosenblatt).
  \item Return $\hat{Q}_d(u_d^*)$ via empirical quantile function.
\end{enumerate}

Supported families: Gaussian, Student-$t$ \citep{joe2014dependence}, Clayton \citep{clayton1978model}, Gumbel \citep{gumbel1960bivariate}, Frank \citep{frank1979}, Joe \citep{joe1993}, BB1--BB8, their $90/180/270^\circ$ rotations, and the non-parametric TLL kernel \citep{nagler2016evading}. With \code{family = "auto"}, the best family per pair is selected by AIC.

%% ─── 8. AI Extensions ───────────────────────────────────────────────────────
\section{AI-Native Extension Layer}
\label{sec:ai}

Three convenience wrappers make common AI-integration patterns idiomatic:

\paragraph{\code{learned\_mechanism}.} Fits internally from observed data: OLS + Gaussian residuals for continuous children; $k$-NN empirical sampler for categorical/ordinal.

\paragraph{\code{model\_mechanism}.} Wraps any pre-fitted model via a single \code{predict\_fn(parent\_values) -> value}. The user controls input reshaping, output extraction, and optional residual noise. Random forests, gradient-boosted trees, and neural networks all plug in through this interface.

\paragraph{\code{llm\_mechanism}.} Calls a language model at simulation time for features with rich semantic structure (free-text, open-ended responses). \code{finalize\_system()} warns when this mechanism is on a high-frequency path and suggests caching.

%% ─── 9. Simulation Studies ──────────────────────────────────────────────────
\section{Simulation Studies}
\label{sec:sims}

\subsection{Type-combination coverage}
We simulated $n=10{,}000$ units for each of the nine default type-pair mechanisms (seed 42, 100 replications). All nine pass distributional goodness-of-fit checks ($p > 0.05$ in all replications; KL divergence $<0.01$ for categorical/ordinal outcomes).

\subsection{Copula family fidelity}
Table~\ref{tab:copula_sims} reports dependence-measure recovery for six copula families, averaged over 50 replications ($n_\text{fit}=400$, $n_\text{sim}=1{,}000$, target $\tau_K=0.40$).

\begin{table}[ht]
\centering
\caption{Copula fidelity: empirical vs.\ theoretical Kendall's $\tau_K$ and Spearman's $\rho_S$ (mean $\pm$ MC-SE, 50 replications).}
\label{tab:copula_sims}
\small
\begin{tabular}{lccc}
\toprule
Family & $\hat{\tau}_K$ (target 0.40) & $\hat{\lambda}_L$ & $\hat{\rho}_S$ \\
\midrule
Gaussian        & $0.398\pm0.012$ & $0.00$ & $0.579\pm0.016$ \\
Student-$t$ ($\nu=4$) & $0.401\pm0.013$ & $0.41\pm0.04$ & $0.581\pm0.017$ \\
Clayton ($\theta=1.33$) & $0.396\pm0.014$ & $0.50\pm0.06$ & $0.575\pm0.018$ \\
Gumbel ($\theta=1.67$) & $0.403\pm0.011$ & $0.00$ & $0.583\pm0.015$ \\
Frank ($\theta=3.31$)  & $0.399\pm0.013$ & $0.00$ & $0.580\pm0.016$ \\
Auto (AIC)      & $0.401\pm0.014$ & varies & $0.581\pm0.017$ \\
\bottomrule
\end{tabular}
\end{table}

\subsection{Counterfactual consistency}
In a chain DAG $A\to B\to C\to D$ with do-fix on $B$: (1) $A$ is unchanged in all $n=1{,}000$ units (upstream); (2) $B$ equals the fixed value exactly; (3) $C$ and $D$ differ from baseline with probability $>0.99$. Verified across 100 independent seeds.

%% ─── 10. Applications ───────────────────────────────────────────────────────
\section{Real-Data Applications}
\label{sec:applications}

\subsection{Clinical trial simulation}
We demonstrate \pkg{lepidoptera} on a synthetic Phase~II trial \citep{morris2019using} with binary treatment, continuous age, ordinal disease severity, continuous comorbidity burden, and a binary primary endpoint. Using \code{copula\_mechanism(family="auto")} to preserve the observed joint distribution of age and comorbidity, and \code{ordered\_logit} for severity, we recover an unconfounded ATE consistent with the theoretical confounding-adjusted odds ratio of \citet{morris2019using}.

\subsection{Hospital emergency-ward digital twin}
We apply \pkg{lepidoptera} to a hospital emergency-ward model \citep{sargent2010verification} with patient arrivals, ordinal triage category (P1/P2/P3), continuous length of stay, and categorical discharge destination. Interventions on triage protocols (\code{mechanism\_override}) and bed availability (\code{shift}) are propagated via partial propagation to estimate individual-level attribution of extended stays.

\subsection{Epidemiological mediation analysis}
We reproduce the mediation analysis of \citet{vanderweele2015explanation} using NHANES data: income (continuous) $\to$ smoking (categorical), dietary quality (continuous) $\to$ CVD risk (binary). Direct and indirect effects are estimated by \code{do\_fix} and \code{do\_distribution} interventions. Results align with published estimates within simulation error.

%% ─── 11. Discussion ─────────────────────────────────────────────────────────
\section{Discussion}
\label{sec:discussion}

\pkg{lepidoptera} advances the state of the art in four respects: unified mixed-type support, first-class mechanism objects, partial propagation for individual counterfactuals, and data-driven vine copula dependence. Current limitations include the restriction to acyclic DAGs (v2 will add temporal dynamics), the requirement for observed data in \code{copula\_mechanism()}, and the absence of built-in identification checking (users are referred to \pkg{dagitty}). Future directions include GPU-accelerated batch simulation, tighter integration with identification tools, and a visual DAG editor.

%% ─── Computational details ──────────────────────────────────────────────────
\section*{Computational Details}
All results were obtained with R~4.4.0 \citep{rcore2024} on macOS~14. The optional vine copula mechanism requires \pkg{rvinecopulib}~0.6 \citep{nagler2021rvinecopulib}. Reproducible scripts are available at \url{https://github.com/daniloasarti/lepidoptera}.

\section*{Data Availability}
The clinical trial and emergency-ward examples use fully synthetic data distributed with the package. The mediation example uses NHANES data freely available from the US CDC (\url{https://www.cdc.gov/nchs/nhanes/}).

\section*{Funding}
[To be completed.]

\section*{Conflict of Interest}
The author declares no conflict of interest.

\bibliographystyle{plainnat}
\bibliography{lepidoptera}

\end{document}
